\documentclass[10pt]{mypackage}

%\usepackage{mlmodern}
%\usepackage{newpxtext,eulerpx,eucal}
%\renewcommand*{\mathbb}[1]{\varmathbb{#1}}

%\usepackage{homework}
\usepackage{notes}

\usepackage[ backend=bibtex, style = alphabetic, sorting=ynt ]{biblatex}
\addbibresource{all_references.bib}

\usepackage{parskip}

\fancyhf{}
\fancyhead[R]{Avinash Iyer}
\fancyhead[L]{Von Neumann Algebras of Equivalence Relations}
\fancyfoot[C]{\thepage}

\setcounter{secnumdepth}{0}

\begin{document}
\RaggedRight
\section{Orbit Equivalence}%
We will assume that $\left( X,\mu \right)$ is a nonatomic standard Borel probability space (i.e., a Borel probability space over a separable completely metrizable space $X$).
\begin{definition}
  A \textit{countable Borel equivalence relation} $\mathcal{R}\subseteq X\times X$ is a Borel subset and equivalence relation satisfying the assumption that
  \begin{align*}
    \left[ x \right]_{\mathcal{R}} &\coloneq \set{y\in X | \left( x,y \right)\in \mathcal{R}} 
  \end{align*}
  is countable.
\end{definition}
The set $\left[ \mathcal{R} \right]$ is known as the \textit{full group} of the equivalence relation $\mathcal{R}$ and consists of all Borel isomorphisms $\phi\colon X\rightarrow X$ with its graph satisfying $\operatorname{graph}\left(\phi\right)\subseteq \mathcal{R}$.

The set of all \textit{partial} Borel isomorphisms $\phi\colon \operatorname{dom}\left( \phi \right)\rightarrow \operatorname{ran}\left( \phi \right)$ with $\operatorname{graph}\left( \phi \right)\subseteq \mathcal{R}$ will be denoted $\left[ \left[ \mathcal{R} \right] \right]$.

The most basic example of a countable Borel equivalence relation is the space of orbits from the action of a countable discrete group $\Gamma\curvearrowright X$. We will write this as $\mathcal{R}\left( \Gamma\curvearrowright X \right)$. From now on, we will use ``equivalence relation'' to refer to countable Borel equivalence relations.

In fact, one can show that all such equivalence relations arise in this form.
\begin{theorem}[Feldman--Moore Theorem]
  If $\mathcal{R}$ is a countable Borel equivalence relation on $X$, then there exists a countable group $\Gamma$ and a Borel action $\Gamma\curvearrowright X$ such that $\mathcal{R} = \mathcal{R}\left( \Gamma\curvearrowright X \right)$. Moreover, both $\Gamma$ and the action can be chosen such that $\left( x,y \right)\in \mathcal{R}$ if and only if there is $g\in \Gamma$ with order $2$ such that $y = g\cdot x$.
\end{theorem}
We say the measure $\mu$ on $X$ is $\mathcal{R}$-invariant if
\begin{align*}
  \phi_{\ast}\mu &= \mu
\end{align*}
for all $\phi\in \left[ \mathcal{R} \right]$, where $\phi_{\ast}\mu\left( U \right) = \mu\left( \phi^{-1}\left( U \right) \right)$ for all Borel $U\subseteq X$.
\end{document}
